\documentclass[11pt]{article}
\usepackage[a4paper,margin=2.4cm]{geometry}
\usepackage{amsmath,amssymb,amsthm}
\usepackage{mathtools}
\usepackage{enumitem}
\usepackage[colorlinks=true,linkcolor=blue,urlcolor=blue]{hyperref}
\newcommand{\R}{\mathbb{R}}
\newcommand{\E}{\mathbb{E}}
\newcommand{\Var}{\operatorname{Var}}
\newcommand{\Cov}{\operatorname{Cov}}
\newcommand{\dd}{\,\mathrm{d}}
\newtheorem{remark}{Remark}

\title{Lecture 04 --- Monte Carlo Simulation for Option Pricing}
\author{Computational Finance (WI4151), TU Delft --- Prof.\ Fang Fang}
\date{}

\begin{document}
\maketitle

\noindent
These notes cover Cluster 2: Monte Carlo (MC) integration, MC simulation of SDEs
(Euler--Maruyama, Milstein), strong/weak convergence orders, and the four variance
reduction techniques (antithetic variates, control variates, quasi-MC, importance
sampling). The guiding theme: MC is the slow but dimension-robust benchmark against
which the fast deterministic methods of this course --- COS (Lectures 6, 10) and
Carr--Madan FFT --- are measured. MC converges at the model-independent rate
$O(N^{-1/2})$; COS converges \emph{exponentially} in the number of terms for smooth
densities. Knowing precisely why MC is $O(N^{-1/2})$ is the way to articulate why COS
is worth the trouble.

\tableofcontents

%==================================================================
\section{Monte Carlo integration and the $O(N^{-1/2})$ error}
%==================================================================

\subsection{The estimator}

Let $I=\int_a^b f(x)\dd x$ be a one-dimensional integral. Write it as an expectation by
multiplying and dividing by the length $b-a$:
\[
I=(b-a)\int_a^b f(x)\,\frac{1}{b-a}\dd x=(b-a)\,\E[f(U)],\qquad U\sim\mathcal U[a,b],
\]
since $\tfrac{1}{b-a}\mathbf 1_{[a,b]}$ is exactly the density of the uniform law on
$[a,b]$. Drawing i.i.d.\ $u_1,\dots,u_N\sim\mathcal U[a,b]$, the estimator is
\[
\widehat I_N=\frac{b-a}{N}\sum_{i=1}^N f(u_i).
\]
The general (multi-dimensional) form factorises the integrand $f(\mathbf x)=h(\mathbf
x)g(\mathbf x)$ with $g$ a probability density on $\Omega\subseteq\R^d$, so that
$I=\int_\Omega h\,g\dd\mathbf x=\E_g[h(X)]$, estimated by
$\widehat I_N=\frac1N\sum_{i=1}^N h(x_i)$ with $x_i\sim g$. This factorisation is the
whole game: \emph{any} integral becomes an average of i.i.d.\ samples once you identify
a density to sample from.

\subsection{Unbiasedness, consistency, and the standard error}

Let $X_i=h(x_i)$, $\mu=\E[h(X)]=I$, $\sigma^2=\Var(h(X))$ (assume $\sigma^2<\infty$).
The estimator $\bar X_N=\frac1N\sum_{i=1}^N X_i$ satisfies:

\paragraph{Unbiasedness.} $\E[\bar X_N]=\frac1N\sum_{i=1}^N\E[X_i]=\mu=I$.

\paragraph{Consistency.} The Strong Law of Large Numbers gives
$\bar X_N\xrightarrow{\text{a.s.}}\mu$, so the estimator converges to the true integral
for almost every realisation of the random stream. (The Weak Law gives the same in
probability; a.s.\ is the stronger statement and is what the slides assert.)

\paragraph{Variance of the estimator.} Using independence,
\[
\Var(\bar X_N)=\Var\!\Bigl(\frac1N\sum_{i=1}^N X_i\Bigr)
=\frac{1}{N^2}\sum_{i=1}^N\Var(X_i)
=\frac{1}{N^2}\cdot N\sigma^2=\frac{\sigma^2}{N}.
\]
The cross terms vanish because $\Cov(X_i,X_j)=0$ for $i\ne j$. Hence the
\emph{standard error} is
\[
\operatorname{se}(\bar X_N)=\sqrt{\Var(\bar X_N)}=\frac{\sigma}{\sqrt N}=O(N^{-1/2}).
\]
This is the central fact of the whole lecture: \textbf{the MC error decays like
$N^{-1/2}$, independent of the dimension $d$}. To gain one decimal digit (factor $10$
in accuracy) you need $100\times$ the samples. By contrast, a tensor-product quadrature
rule of order $k$ in $d$ dimensions costs $N=m^d$ points for $m$ points per axis and
achieves $O(m^{-k})=O(N^{-k/d})$ --- the rate \emph{degrades with $d$} (curse of
dimensionality). MC's rate does not. This is why MC dominates in high dimension (basket
options, many time steps) despite its slowness in low dimension.

\subsection{Confidence intervals from the CLT}

The Central Limit Theorem sharpens the standard error into a probabilistic error bar.
With $\mu=I$ and $\sigma^2=\Var(X)$,
\[
\frac{\bar X_N-I}{\sigma/\sqrt N}\xrightarrow{d}\mathcal N(0,1).
\]
Therefore, for large $N$, $\bar X_N\approx\mathcal N(I,\sigma^2/N)$, and an
asymptotic $(1-\alpha)$ confidence interval for $I$ is
\[
\Bigl[\,\bar X_N-z_{1-\alpha/2}\,\frac{\sigma}{\sqrt N},\;\;
\bar X_N+z_{1-\alpha/2}\,\frac{\sigma}{\sqrt N}\,\Bigr],
\]
where $z_{1-\alpha/2}$ is the standard normal quantile (e.g.\ $1.96$ for $95\%$). In
practice $\sigma$ is unknown and is replaced by the sample standard deviation
\[
s_N^2=\frac{1}{N-1}\sum_{i=1}^N(X_i-\bar X_N)^2,
\]
(the $N-1$ gives an unbiased variance estimate). The half-width
$z_{1-\alpha/2}\,s_N/\sqrt N$ is the quantity actually reported alongside an MC price; a
variance-reduction technique is "good" precisely when it shrinks this half-width for
fixed $N$ (equivalently, the same accuracy at fewer samples).

\begin{remark}
The three convergence modes in the slides form a strict hierarchy:
$\xrightarrow{\text{a.s.}}\Rightarrow\xrightarrow{p}\Rightarrow\xrightarrow{d}$, with
no reverse implications in general. The SLLN is a.s., the WLLN is in probability, the
CLT is in distribution. The CLT does \emph{not} say $\bar X_N\to I$ in distribution
(that would be degenerate); it says the \emph{rescaled} error
$\sqrt N(\bar X_N-I)/\sigma$ is asymptotically standard normal.
\end{remark}

%==================================================================
\section{Pseudo-random vs.\ quasi-random (QMC)}
%==================================================================

\subsection{Why QMC can beat $O(N^{-1/2})$}

Standard MC uses pseudo-random points that mimic i.i.d.\ uniforms; its error is
\emph{random} and decays at $O(N^{-1/2})$ by the CLT. Quasi-Monte Carlo (QMC) replaces
the random points by a \emph{deterministic} low-discrepancy sequence
$u_1,\dots,u_N\in[0,1]^d$ (Sobol, Halton, Faure, lattice rules) that fills the cube more
evenly than random points, which clump and leave gaps. Both methods compute the
\emph{same} estimator $\frac1N\sum_i f(x_i)$; they differ only in how the $x_i$ are
chosen.

The error is governed not by the CLT but by the deterministic \textbf{Koksma--Hlawka
inequality}:
\[
\Bigl|\frac1N\sum_{i=1}^N f(x_i)-\int_{[0,1]^d}f(x)\dd x\Bigr|\le V(f)\,D_N^\ast,
\]
a product of two factors that separate the integrand from the point set:
\begin{itemize}[leftmargin=2em]
\item $V(f)$ is the variation of $f$ in the sense of Hardy--Krause, measuring the
total magnitude of all mixed partial derivatives across coordinate subsets --- how
"rough" $f$ is. It depends only on the integrand.
\item $D_N^\ast$ is the \emph{star discrepancy} of the point set
$P_N=\{x_1,\dots,x_N\}$,
\[
D_N^\ast(P_N)=\sup_{t\in[0,1]^d}\Bigl|\frac1N\sum_{i=1}^N
\mathbf 1_{\{x_i\in[0,t)\}}-\prod_{j=1}^d t_j\Bigr|,
\]
the worst-case gap, over all origin-anchored boxes $[0,t)=[0,t_1)\times\cdots\times
[0,t_d)$, between the empirical fraction of points in the box and its Lebesgue volume
$\prod_j t_j$. It depends only on the points.
\end{itemize}
For a low-discrepancy sequence one has
\[
D_N^\ast=O\!\Bigl(\frac{(\log N)^d}{N}\Bigr),
\]
so for sufficiently smooth (bounded-variation) $f$ the QMC error is
$O((\log N)^d/N)$. For fixed $d$ this is \emph{essentially $O(N^{-1})$} up to the
logarithmic factor --- a full power of $N$ faster than the $O(N^{-1/2})$ of plain MC.

\subsection{Caveats}

\begin{itemize}[leftmargin=2em]
\item The rate $(\log N)^d/N$ carries the dimension in the exponent of the logarithm.
For large $d$, $(\log N)^d$ is enormous and the asymptotic advantage only appears at
impractically large $N$. In practice the \emph{effective dimension} (how much each
coordinate actually matters) is what governs QMC performance; Brownian-bridge or PCA
constructions of the path lower the effective dimension and are what make QMC work for
option pricing.
\item QMC is biased-free but \emph{deterministic}, so a plain QMC run gives no error
estimate. \emph{Randomised QMC} (e.g.\ random shifts of a lattice rule,
$x_n=\{nz/N+\Delta\}\bmod 1$ with random $\Delta$) restores an unbiased estimator and
lets you estimate the error from independent shifts, while keeping most of the
discrepancy advantage.
\item Mapping to non-unit / unbounded domains: scale $[0,1]^d$ linearly to
$[a,b]^d$ via $x^{\text{new}}_i=a+(b-a)x_i$, and map to a target law by the inverse CDF,
$x_i=\Phi^{-1}(u_i)$ (for Gaussians). The inverse-CDF transform preserves
low-discrepancy structure better than Box--Muller.
\end{itemize}

\begin{remark}[Connection to COS/FFT]
QMC, like COS, exploits \emph{smoothness} to beat $O(N^{-1/2})$ --- but QMC trades on
smoothness/low variation of the integrand whereas COS trades on smoothness of the
\emph{density}, which it never touches directly, working instead with its Fourier-cosine
coefficients (the characteristic function). When the characteristic function is known in
closed form (Heston, Lévy models), COS achieves exponential convergence, far beyond
even QMC's near-$O(N^{-1})$.
\end{remark}

%==================================================================
\section{Discretising SDEs}
%==================================================================

When the terminal law of $X_T$ is not available in closed form, MC pricing requires
simulating paths of the SDE
\[
\dd X_t=\mu(t,X_t)\dd t+\sigma(t,X_t)\dd W_t.
\]
On a uniform grid $t_k=kh$, $k=0,\dots,K$, $Kh=T$, with i.i.d.\ $Z_k\sim\mathcal N(0,1)$,
the increments $W_{t_k}-W_{t_{k-1}}\stackrel{d}{=}\sqrt h\,Z_k$.

\subsection{Euler--Maruyama}

The simplest scheme freezes the coefficients at the left endpoint of each interval:
\[
\widehat X_{t_k}=\widehat X_{t_{k-1}}+\mu(t_{k-1},\widehat X_{t_{k-1}})\,h
+\sigma(t_{k-1},\widehat X_{t_{k-1}})\,\sqrt h\,Z_k,
\qquad \widehat X_0=X_0.
\]
This is the Itô (left-endpoint) discretisation of
$X_{t_k}=X_{t_{k-1}}+\int_{t_{k-1}}^{t_k}\mu\dd t+\int_{t_{k-1}}^{t_k}\sigma\dd W$,
with each integral approximated by its left-point value times the increment. The
left-endpoint choice is not cosmetic: it is what makes the stochastic integral a
martingale (the Itô convention), as the $\int W\dd W$ example in the slides shows.

\subsection{Milstein: deriving the extra term from the Itô--Taylor expansion}

Euler ignores the variation of $\sigma$ \emph{within} a step. Milstein restores the
leading correction. Work on one step $[t_{k-1},t_k]$, write $t_{k-1}=s$, freeze
$X_s=x$, and consider (for clarity) the autonomous scalar case
$\dd X_t=\mu(X_t)\dd t+\sigma(X_t)\dd W_t$. The exact increment is
\[
X_{t_k}-X_s=\int_s^{t_k}\mu(X_u)\dd u+\int_s^{t_k}\sigma(X_u)\dd W_u.
\]
Euler replaces $\mu(X_u)\approx\mu(x)$ and $\sigma(X_u)\approx\sigma(x)$. The
$\dd t$ integral is already $O(h)$ and its error is higher order, so the gain must come
from the diffusion integral. Expand $\sigma(X_u)$ around $X_s=x$ by Itô's lemma applied
to the function $\sigma$:
\[
\dd\,\sigma(X_t)=\sigma'(X_t)\dd X_t+\tfrac12\sigma''(X_t)(\dd X_t)^2
=\Bigl[\sigma'\mu+\tfrac12\sigma''\sigma^2\Bigr]\dd t+\sigma'\sigma\,\dd W_t,
\]
using the Itô multiplication rule $(\dd W_t)^2=\dd t$. Integrating from $s$ to $u$,
\[
\sigma(X_u)=\sigma(x)+\int_s^u\sigma'(X_r)\sigma(X_r)\dd W_r+\big(\text{$O(\dd t)$ terms}\big).
\]
Substitute the \emph{leading} stochastic correction back into the diffusion integral and
freeze the inner coefficient at $x$ (the rest is higher order):
\[
\int_s^{t_k}\sigma(X_u)\dd W_u
\approx\sigma(x)\,(W_{t_k}-W_s)
+\sigma'(x)\sigma(x)\int_s^{t_k}\!\!\int_s^u \dd W_r\,\dd W_u.
\]
The iterated Itô integral evaluates in closed form. With $\Delta W:=W_{t_k}-W_s$,
\[
\int_s^{t_k}\!\!\int_s^u \dd W_r\,\dd W_u
=\int_s^{t_k}(W_u-W_s)\dd W_u
=\tfrac12\bigl[(\Delta W)^2-(t_k-s)\bigr]
=\tfrac12\bigl[(\Delta W)^2-h\bigr],
\]
where the $-\tfrac12 h$ is exactly the Itô correction (the $\int W\,\dd W=\tfrac12(W^2-t)$
identity from the martingale example). Writing $\Delta W=\sqrt h\,Z_k$, so
$(\Delta W)^2-h=h(Z_k^2-1)$, the one-step update becomes
\[
\boxed{\;\widehat X_{t_k}=\widehat X_{t_{k-1}}+\mu\,h+\sigma\,\sqrt h\,Z_k
+\tfrac12\,\sigma\,\frac{\partial\sigma}{\partial x}\,(Z_k^2-1)\,h\;}
\]
with all coefficients evaluated at $(t_{k-1},\widehat X_{t_{k-1}})$. The extra term
$\tfrac12\sigma\,\partial_x\sigma\,(Z_k^2-1)h$ is the Milstein correction: it captures
the curvature of the diffusion ($\partial_x\sigma\ne0$) over the step. Note it vanishes
for \emph{additive} noise ($\sigma$ constant), where Euler and Milstein coincide --- and
indeed for additive noise Euler already has strong order $1$.

\begin{remark}
Milstein needs the derivative $\partial_x\sigma$ in closed form, and the
multidimensional version requires iterated integrals
$\int\!\int\dd W^{(i)}\dd W^{(j)}$ (Lévy areas) for $i\ne j$ that cannot be written in
terms of the increments alone --- this is why the slides note that Milstein "is not
trivial to generalise to multi-dimensional SDEs." For the Heston variance process
$\dd v_t=\kappa(\bar v-v_t)\dd t+\gamma\sqrt{v_t}\dd W_t$, $\sigma(v)=\gamma\sqrt v$,
$\partial_v\sigma=\gamma/(2\sqrt v)$, so $\sigma\,\partial_v\sigma=\gamma^2/2$ and the
Milstein term is $\tfrac14\gamma^2(Z^2-1)h$ --- this is the standard full-truncation /
Milstein scheme that mitigates (but does not eliminate) the negative-variance problem of
plain Euler on Heston.
\end{remark}

%==================================================================
\section{Strong vs.\ weak convergence}
%==================================================================

There are two error measures for a discretisation scheme with step $h$, depending on
whether we care about the pathwise value or only a distributional functional.

\paragraph{Strong order $\beta$.} Pathwise (root-mean / $L^1$) accuracy:
\[
\E\bigl[\,|\widehat X_{Kh}-X_T|\,\bigr]\le c\,h^\beta
\]
for some constant $c$ and all small $h$. This requires the simulated path and the true
path (driven by the \emph{same} Brownian motion) to stay close.

\paragraph{Weak order $\beta$.} Distributional accuracy: for all sufficiently smooth
test functions $f$ (formally $f\in C_P^{2\beta+2}$, derivatives up to order $2\beta+2$
of polynomial growth),
\[
\bigl|\E[f(\widehat X_{Kh})]-\E[f(X_T)]\bigr|\le c\,h^\beta.
\]
Only the \emph{laws} of $\widehat X_{Kh}$ and $X_T$ need to match, not the paths.

\paragraph{Orders.}
\[
\text{Euler--Maruyama: strong }\beta=\tfrac12,\ \text{weak }\beta=1;\qquad
\text{Milstein: strong }\beta=1,\ \text{weak }\beta=1.
\]
Milstein's payoff is the \emph{strong} rate ($\tfrac12\to1$); its weak rate equals
Euler's. So if all you want is a European price --- an expectation $\E[f(X_T)]$ --- Euler
already gives weak order $1$ and Milstein buys you nothing distributionally. Milstein
matters for \emph{pathwise} accuracy: path-dependent payoffs (barriers, Asians,
lookbacks) where the whole trajectory enters, and for any scheme where the same noise is
reused (e.g.\ multilevel MC, which exploits strong convergence to couple coarse and fine
levels).

\begin{remark}[The $h^{1/3}$ reality check]
The weak-order-$1$ statement requires the test function $f$ to be smooth. Option
payoffs like $(S_T-K)^+$ are \emph{not} differentiable at the strike, so the theoretical
$\beta=1$ is not guaranteed. Broadie--Kaya (2006) report empirical rates of
$O(h^{1/3})$ or worse for Euler on European prices under stochastic volatility. The
slides' heuristic: balance bias $\sim h$ against MC noise $\sim N^{-1/2}$ at total cost
$s\sim N/h$. Eliminating $N=sh$ gives error $\sim h+ (sh)^{-1/2}$; optimising over $h$
(setting $h\sim(sh)^{-1/2}$, i.e.\ $h\sim s^{-1/3}$) yields overall error $O(s^{-1/3})$
in the computational budget $s$. This is much worse than the deterministic methods of
this course and is the practical motivation for COS/FFT on models with a known
characteristic function: they sidestep path discretisation entirely.
\end{remark}

%==================================================================
\section{Variance reduction}
%==================================================================

All four techniques attack the same target: shrink $\Var(\bar X_N)=\sigma^2/N$ by
lowering the effective $\sigma^2$, keeping the estimator unbiased (or correcting the
bias). A method giving variance $\rho\,\sigma^2$ with $\rho<1$ delivers the accuracy of
$N/\rho$ plain samples at cost $N$ --- a speed-up of $1/\rho$.

\subsection{Antithetic variates}

\paragraph{Construction (uniform case).} To estimate $I=\E[f(U)]$, $U\sim\mathcal
U(0,1)$, pair each draw $U_i$ with its reflection $1-U_i$ (also $\mathcal U(0,1)$):
\[
\widehat I_M=\frac1M\sum_{i=1}^M\frac{f(U_i)+f(1-U_i)}{2}.
\]
Since $1-U_i\sim\mathcal U(0,1)$, each summand has mean $I$, so $\E[\widehat I_M]=I$:
the estimator is unbiased.

\paragraph{Variance.} For one pair,
\[
\Var\!\Bigl(\frac{f(U)+f(1-U)}2\Bigr)
=\tfrac14\Var(f(U))+\tfrac14\Var(f(1-U))+\tfrac12\Cov(f(U),f(1-U)).
\]
Because $U$ and $1-U$ have the same law, $\Var(f(U))=\Var(f(1-U))$, so this collapses to
\[
\Var\!\Bigl(\frac{f(U)+f(1-U)}2\Bigr)
=\tfrac12\Var(f(U))+\tfrac12\Cov(f(U),f(1-U)),
\]
hence $\Var(\widehat I_M)=\frac1M\bigl[\tfrac12\Var(f(U))+\tfrac12\Cov(f(U),f(1-U))\bigr]$.

\paragraph{When it helps.} Compare with the variance of the \emph{plain} estimator
using the same $2M$ function evaluations, $\frac{1}{2M}\Var(f(U))$. Antithetic wins iff
\[
\Cov(f(U),f(1-U))<0.
\]
This is guaranteed when $f$ is monotone. The argument: if $f$ is monotone increasing,
then $g(x):=-f(1-x)$ is also monotone increasing; by the FKG/correlation lemma (two
comonotone functions of the same random variable have nonnegative covariance),
$\Cov(f(U),-f(1-U))\ge0$, i.e.\ $\Cov(f(U),f(1-U))\le0$. (The lemma: if $f,g$ are both
increasing, $\Cov(f(X),g(X))\ge0$.) The Gaussian version is identical with the
reflection $U\mapsto-U$ (both $\mathcal N(0,1)$): along a path one runs the mirror path
$-Z_0,\dots,-Z_{N-1}$ and averages the two payoffs. For a vanilla call the discounted
payoff is monotone in the driving normal, so antithetic always reduces variance there.

\begin{remark}
Antithetic can \emph{fail} (raise variance) when $f$ is strongly non-monotone, e.g.\ a
symmetric straddle or a payoff symmetric about the forward, where $f(U)$ and $f(1-U)$
become positively correlated. The monotonicity condition is the safe guarantee.
\end{remark}

\subsection{Control variates: the optimal coefficient $c^\star$}

\paragraph{Setup.} We want $\E[X]$ (e.g.\ the arithmetic-Asian price). Suppose we can
simulate, on the same paths, a correlated $Y$ (e.g.\ the geometric-Asian payoff) whose
mean $\E[Y]$ is \emph{known} in closed form. For any coefficient $\theta$ define
\[
Z_\theta=X+\theta\bigl(\E[Y]-Y\bigr).
\]
Then $\E[Z_\theta]=\E[X]+\theta(\E[Y]-\E[Y])=\E[X]$ for every $\theta$: unbiased
regardless of $\theta$. Only the variance depends on $\theta$:
\[
\Var(Z_\theta)=\Var(X)-2\theta\Cov(X,Y)+\theta^2\Var(Y).
\]

\paragraph{Optimisation.} This is a quadratic in $\theta$, minimised by setting the
derivative $-2\Cov(X,Y)+2\theta\Var(Y)=0$:
\[
\boxed{\;\theta^\star=\frac{\Cov(X,Y)}{\Var(Y)}\;}
\]
(Completing the square,
$\Var(Z_\theta)=\Var(Y)\bigl(\theta-\tfrac{\Cov(X,Y)}{\Var(Y)}\bigr)^2
+\Var(X)-\tfrac{\Cov(X,Y)^2}{\Var(Y)}$, exposing both the minimiser and the minimum.)
Note: with the slides' sign convention $Z=X+\theta(\E Y-Y)$, the optimum is
$\theta^\star=+\Cov(X,Y)/\Var(Y)$; if instead one writes $Z=X+c(Y-\E Y)$, the same
algebra gives $c^\star=-\Cov(X,Y)/\Var(Y)$ --- the textbook "$c^\star=-\Cov/\Var$." They
are the same estimator; only the sign convention differs.

\paragraph{Resulting variance reduction.} Substituting $\theta^\star$,
\[
\Var(Z_{\theta^\star})=\Var(X)-\frac{\Cov(X,Y)^2}{\Var(Y)}
=\bigl(1-\rho_{XY}^2\bigr)\Var(X),
\qquad
\rho_{XY}=\frac{\Cov(X,Y)}{\sqrt{\Var(X)\Var(Y)}}.
\]
The variance is reduced by the factor $1-\rho_{XY}^2$. The reduction is dramatic exactly
when $X$ and $Y$ are highly correlated: a control variate with $|\rho|=0.99$ cuts
variance by $\approx 50\times$. The classic example: price the arithmetic-average Asian
call (no formula) using the geometric-average Asian call (closed form for equally spaced
fixings) as $Y$; the two payoffs are nearly perfectly correlated along each path.

\paragraph{Practical point.} $\Cov(X,Y)$ and $\Var(Y)$ are usually unknown and are
\emph{estimated from the same simulation} (a regression of $X$ on $Y$ gives
$\widehat\theta^\star$). Plugging in an estimated $\widehat\theta^\star$ introduces a
small $O(1/N)$ bias, negligible against the $O(N^{-1/2})$ statistical error; a clean fix
is to estimate $\theta^\star$ on a pilot run and apply it on an independent run.

\subsection{Importance sampling}

\paragraph{Idea and the likelihood ratio.} To estimate
$\E_p[f(X)]=\int f(x)p(x)\dd x$ when the "important" region (where $f\,p$ is large) is
rarely hit by $p$ --- the canonical case being deep out-of-the-money options, where the
payoff is nonzero only on a low-probability set --- sample instead from a proposal $q$
and reweight:
\[
\E_p[f(X)]=\int f(x)\,\frac{p(x)}{q(x)}\,q(x)\dd x
=\E_q\!\Bigl[f(X)\,\frac{p(X)}{q(X)}\Bigr].
\]
The factor $p(X)/q(X)$ is the \emph{likelihood ratio} (Radon--Nikodym derivative
$\dd p/\dd q$); it corrects the bias from sampling under the "wrong" measure, leaving the
estimator unbiased. The estimator is
$\frac1N\sum_i f(X_i)\,\tfrac{p(X_i)}{q(X_i)}$ with $X_i\sim q$.

\paragraph{Gaussian drift shift.} For $V=\E[f(X)]$ with $X\sim\mathcal N(0,I_d)$, the
natural proposal is a mean shift $X=Z+\theta$, $Z\sim\mathcal N(0,I_d)$. The
likelihood ratio of $\mathcal N(0,I)$ w.r.t.\ $\mathcal N(\theta,I)$ evaluated at the
sampled point is $\exp(-\theta^\top Z-\tfrac12\|\theta\|^2)$, so
\[
V=\E\Bigl[f(Z+\theta)\,\exp\bigl(-\theta^\top Z-\tfrac12\|\theta\|^2\bigr)\Bigr].
\]
(Scalar derivation: with $f_{0}(z)=\tfrac1{\sqrt{2\pi}}e^{-z^2/2}$ and
$g_\mu(z)=\tfrac1{\sqrt{2\pi}}e^{-(z-\mu)^2/2}$, the ratio is
$f_0(z)/g_\mu(z)=\exp(-\tfrac{z^2}2+\tfrac{(z-\mu)^2}2)=\exp(-\mu z+\tfrac12\mu^2)$.) The
shift $\theta$ is chosen to move mass into the payoff region (e.g.\ for a call, shift the
driving normal so that $S_T>K$ becomes typical). Optimal $\theta^\star$ minimises the
second moment
$M(\theta)=\E[f(Z+\theta)^2\exp(-2\theta^\top Z-\|\theta\|^2)]$, solving
$\nabla_\theta M(\theta^\star)=0$.

\paragraph{The double-edged sword.} A well-chosen $q$ can reduce variance by orders of
magnitude; a poorly chosen one (light tails relative to $f\,p$, so the ratio
$p/q$ blows up) gives an estimator with \emph{infinite} variance --- worse than plain
MC. Importance sampling is the only one of the four techniques that can backfire
catastrophically. Choosing $q$ is "the art of importance sampling."

\begin{remark}[Girsanov / change of numeraire]
In continuous time the change of measure is the Girsanov theorem: shifting the drift of
$W$ by $Y_s$ corresponds to the density process $Z_t^{P\to Q}=\exp(\int_0^t Y_s\dd
W_s^P-\tfrac12\int_0^t Y_s^2\dd s)$, a martingale (under Novikov) with $\E^P[Z_t]=1$.
\emph{Change of numeraire} is a special, financially-motivated change of measure: using
the stock $S_t$ as numeraire turns $\E^Q[S_T\mathbf 1_{\{S_T>K\}}]$ into $e^{rT}S_0
Q^S(S_T>K)=e^{rT}S_0 N(d_1)$, recovering the $N(d_1)$ term of Black--Scholes by a measure
change rather than by integration. Importance sampling does not require the proposal to
come from a numeraire, but the two share the Radon--Nikodym machinery.
\end{remark}

%==================================================================
\section{Where this sits in the course}
%==================================================================

\begin{itemize}[leftmargin=2em]
\item \textbf{MC is the benchmark.} Every fast method in this course --- COS (L6, L10),
Carr--Madan FFT, finite differences (L3, L5), trees (L1) --- is validated against a
brute-force MC price with a CLT confidence interval. When the oral asks "how do you know
your COS price is right," the honest answer is: it agrees with MC to within the MC
standard error, and MC is unbiased by construction.
\item \textbf{Rate comparison.} MC: $O(N^{-1/2})$, dimension-free. QMC:
$O((\log N)^d/N)$, near-$O(N^{-1})$ in low effective dimension. Quadrature: $O(N^{-k/d})$,
cursed in $d$. COS: \emph{exponential} in the number of cosine terms for densities with
known, smooth characteristic functions. This ordering is exactly why the course pivots
to Fourier methods: for a 1-D European option under an affine model, COS reaches machine
precision with a few hundred terms, where MC would need $\sim10^{12}$ paths.
\item \textbf{When MC still wins.} High-dimensional baskets, path-dependent payoffs with
no transform structure, and American/Bermudan via Longstaff--Schwartz (L9) --- MC's
dimension-robustness keeps it indispensable. Variance reduction is what makes it
competitive in those regimes.
\end{itemize}

%==================================================================
\section*{Honest gaps / what the slides do \emph{not} prove}
%==================================================================
\addcontentsline{toc}{section}{Honest gaps}

\begin{itemize}[leftmargin=2em]
\item \textbf{The slides do not derive Milstein from the Itô--Taylor expansion.} They
state the scheme and the convergence orders but only \emph{assert} that the extra term
"accounts for the curvature of the diffusion." The full derivation (Section~3.2 above) is
mine, filling that gap; an examiner could reasonably ask for it.
\item \textbf{The strong/weak orders are stated, not proved.} The slides give the
definitions and quote $\beta_{\text{Euler}}=(\tfrac12,1)$,
$\beta_{\text{Milstein}}=(1,1)$ without proof. Proofs (Kloeden--Platen) are out of scope
but worth knowing exist.
\item \textbf{The optimal importance-sampling drift is left implicit.} The slides reduce
to the first-order condition $\nabla_\theta M(\theta^\star)=0$ but do not solve it in
closed form (there is none in general; for a call one uses the asymptotically optimal
shift placing the mode of the shifted integrand at the strike).
\item \textbf{No explicit confidence-interval slide.} The CLT and standard error are on
the slides, but the confidence-interval formula and the sample-variance estimator
$s_N^2$ (Section~1.3) are standard background I have added.
\item \textbf{QMC's effective-dimension story is absent.} The slides give the
Koksma--Hlawka bound and the $(\log N)^d/N$ rate but do not discuss why QMC works in
practice for high-dimensional path simulation (Brownian bridge / PCA path construction
to lower effective dimension).
\end{itemize}

\end{document}
